\section*{Parte 2: Selección única}

Para esta parte, debe entregar un archivo \texttt{parte2\_NombreApellido.txt} con las respuestas en el siguiente formato:

\begin{minted}{text}
1. X
2. X
3. X
4. X
5. X
\end{minted}

Cada pregunta vale \textbf{6 puntos}. Las preguntas se basan en el siguiente programa, que implementa un codificador y decodificador de código Morse:

\begin{minted}{python}
MORSE = {
    "A": ".-",   "B": "-...", "C": "-.-.", "D": "-..",
    "E": ".",    "F": "..-.", "G": "--.",  "H": "....",
    "I": "..",   "J": ".---", "K": "-.-",  "L": ".-..",
    "M": "--",   "N": "-.",   "O": "---",  "P": ".--.",
    "Q": "--.-", "R": ".-.",  "S": "...",  "T": "-",
    "U": "..-",  "V": "...-", "W": ".--",  "X": "-..-",
    "Y": "-.--", "Z": "--..",
}

def codificar(mensaje):
    resultado = ""
    for letra in mensaje.upper():
        if letra == " ":
            resultado += "/ "
        elif letra in MORSE:
            resultado += MORSE[letra] + " "
    return resultado.strip()

def decodificar(codigo):
    inverso = {}
    for letra, simbolo in MORSE.items():
        inverso[simbolo] = letra
    resultado = ""

    # .split(" / "): Separa el string en una lista según " / "
    for palabra in codigo.split(" / "): 
        for simbolo in palabra.split():
            resultado += inverso[simbolo]
        resultado += " "
    return resultado.strip()

mensaje = input("Mensaje: ")
codificado = codificar(mensaje)
print("Morse:", codificado)
print("Decodificado:", decodificar(codificado))
\end{minted}

\subsubsection*{Pregunta 1}

¿Qué verifica la condición \mintinline{python}{elif letra in MORSE} dentro de \texttt{codificar}?

\begin{enumerate}[label=\alph*)]
    \item Que \mintinline{python}{letra} sea un símbolo Morse como \mintinline{python}{".-"} o \mintinline{python}{"-..."}
    \item Que \mintinline{python}{letra} esté definida como clave del diccionario, como \mintinline{python}{"A"} o \mintinline{python}{"B"}
    \item Que \mintinline{python}{letra} sea un carácter alfabético en minúsculas
    \item Que \mintinline{python}{letra} no sea un espacio ni un número
\end{enumerate}

\subsubsection*{Pregunta 2}

¿Qué retorna \mintinline{python}{codificar("Hi!")}?

\begin{enumerate}[label=\alph*)]
    \item \mintinline{text}{".... .. !"}
    \item Un error porque \mintinline{python}{"!"} no existe en el diccionario
    \item \mintinline{text}{".... .. "}
    \item \mintinline{text}{".... .."}
\end{enumerate}

\subsubsection*{Pregunta 3}

¿Cuál es el valor de \mintinline{python}{codificado} después de llamar \mintinline{python}{codificar("ET")}?

\begin{enumerate}[label=\alph*)]
    \item \mintinline{text}{".-"}
    \item \mintinline{text}{"E T"}
    \item \mintinline{text}{". -"}
    \item \mintinline{text}{"E. T-"}
\end{enumerate}

\subsubsection*{Pregunta 4}

¿Cuántas veces se ejecuta la línea \mintinline{python}{resultado += " "} dentro de \texttt{decodificar} si recibe el string \mintinline{text}{".... .-.. / .-"}?

\begin{enumerate}[label=\alph*)]
    \item 1 vez
    \item 3 veces
    \item 4 veces
    \item 2 veces
\end{enumerate}

\subsubsection*{Pregunta 5}

¿Qué retorna \mintinline{python}{decodificar(codificar("hola"))}?

\begin{enumerate}[label=\alph*)]
    \item \mintinline{python}{"hola"}
    \item \mintinline{python}{".... --- .-.. .-"}
    \item \mintinline{python}{"HOLA"}
    \item Un error porque \mintinline{python}{codificar} no acepta minúsculas
\end{enumerate}
